\section{Resonant normal-form extraction}
\label{app:normal-form}

\subsection{Symplectic convention}

In normal-form coordinates \(w=(Q,P)\), use
\[
\boxed{
\omega=dQ\wedge dP,\qquad
\iota_{X_H}\omega=dH.
}
\]
Thus
\[
\boxed{
X_H=(H_P,-H_Q).
}
\]

\subsection{Critical return jet}

At resonance,
\[
P_\star(w)
=
-w+
P_2(w)+P_3(w)+P_4(w)+P_5(w)+\bigO(|w|^6).
\]

The nonlinear gauge is fixed by using only the nonresonant canonical generators required to eliminate the even map terms.  Let $G_3$ be homogeneous of degree three and $Y_2=X_{G_3}$.  Since the homological operator at the linear map $-I$ multiplies an even vector field by two, the quadratic term is removed by
\[
Y_2=-\frac12P_2.
\]
We conjugate by the time-one Hamiltonian flow $\exp(Y_2)$.  If $P_4^{(1)}$ denotes the resulting quartic map term, a homogeneous quintic generator $G_5$ with
\[
Y_4=X_{G_5}=-\frac12P_4^{(1)}
\]
removes it.  No resonant quartic Hamiltonian generator is added; this condition fixes the sixth-order gauge used for the coefficients below.  The resulting map is
\[
\boxed{
\widetilde P(w)
=
-w+\widetilde P_3(w)+\widetilde P_5(w)+\bigO(|w|^7).
}
\]

Define
\[
\boxed{
M=-\widetilde P.
}
\]
Then
\[
M=I+M_3+M_5+\bigO(7).
\]

\subsection{Hamiltonian interpolation}

Seek
\[
M=\exp(X_H),
\qquad
H=H_4+H_6+\bigO(8).
\]
Comparison of homogeneous orders gives
\[
\boxed{
X_{H_4}=M_3,
}
\]
and
\[
\boxed{
X_{H_6}
=
M_5-\frac12DM_3M_3.
}
\]

Let
\[
R_5=M_5-\frac12DM_3M_3.
\]
Then
\[
\boxed{
H_{4,P}=M_3^{(Q)},\qquad
H_{4,Q}=-M_3^{(P)},
}
\]
and
\[
\boxed{
H_{6,P}=R_5^{(Q)},\qquad
H_{6,Q}=-R_5^{(P)}.
}
\]

Hamiltonian consistency requires
\[
\partial_QM_3^{(Q)}+\partial_PM_3^{(P)}=0,
\]
and
\[
\partial_QR_5^{(Q)}+\partial_PR_5^{(P)}=0.
\]

\subsection{Linear reversible normalization}

Away from resonance,
\[
M_\mu=-P_\mu.
\]
The intrinsic detuning is oriented so that
\[
\boxed{
DM_\mu(0)
=
I+
\delta
\begin{pmatrix}
0&-1\\
1&0
\end{pmatrix}
+\bigO(\delta^2).
}
\]
The quadratic generator is
\[
\boxed{
H_2=-\delta I,\qquad I=\frac{Q^2+P^2}{2}.
}
\]

The anisotropic symplectic scale is fixed by
\[
q=\sqrt{\alpha_{\rm lin}}\,Q,\qquad
p=\frac{P}{\sqrt{\alpha_{\rm lin}}}.
\]
The residual rotation commuting with the linear generator is fixed by aligning the \(Q\)- and \(P\)-axes with the tangent fixed sets of the reversing involutions.

\subsection{Independent quartic-balance test}

Before the intrinsic scaling, write the pure quartic terms as
\[
a_4q^4+c_4p^4.
\]
Under
\[
q=\sqrt{\alpha}\,Q,\qquad
p=\frac{P}{\sqrt{\alpha}},
\]
their coefficients become
\[
a_4\alpha^2,\qquad c_4\alpha^{-2}.
\]
Balancing them yields
\[
\boxed{
\alpha_{\rm bal}
=
\left(\frac{c_4}{a_4}\right)^{1/4}.
}
\]

\subsection{Quartic Hamiltonian}

In the fixed reversible chart,
\[
\boxed{
H_4=A(Q^4+P^4)+BQ^2P^2.
}
\]
The independent multiprecision jet transport gives
\[
A=203.6723580,\qquad B=-408.3087120.
\]
Define
\[
\boxed{
\rho^2=2-\frac BA,
}
\qquad
\boxed{
\sigma=4-\rho^2.
}
\]
At \(\rho=2\),
\[
\boxed{
H_4=A(Q^2-P^2)^2.
}
\]

In polar coordinates
\[
Q=\sqrt{2I}\cos\theta,\qquad
P=\sqrt{2I}\sin\theta,
\]
\[
\boxed{
H_4
=
4AI^2-A\rho^2I^2\sin^22\theta.
}
\]
Near \(\theta=\pi/4+x\),
\[
\boxed{
H_4
=
A\sigma I^2+
4A\rho^2I^2x^2+
\bigO(I^2x^4).
}
\]

\subsection{Sixth-order Hamiltonian}

The extracted sixth-order term is
\[
\boxed{
H_6
=
dQ^6+eQ^4P^2+fQ^2P^4+gP^6.
}
\]
In the same nonresonant gauge,
\[
d=27009.26240,\quad e=271555.83499,
\]
\[
f=512441.22879,\quad g=75140.77917.
\]
Define
\[
\boxed{
G=d+e+f+g,
}
\]
and
\[
\boxed{
\Delta_6=(d+e)-(f+g).
}
\]

\subsection{Full detuned normal form}

The truncated generator is
\[
\boxed{
K_\delta
=
-\delta I+H_4+H_6+
\bigO(\delta^2I,\delta I^2,I^4).
}
\]

Since \(F=P^2=M^2\),
\[
\boxed{
F=\exp(2X_{K_\delta})
}
\]
to the retained formal order. Therefore physical second-return exponents and angles are twice those generated by \(M\):
\[
\boxed{
\kappa_F=2\lambda_M,\qquad
\vartheta_F=2\omega_M.
}
\]

\subsection{Frozen-jet axial Floquet expansions}

Within the frozen truncation $-\delta I+H_4+H_6$, for the \(Q\)-axis,
\[
\boxed{
\kappa_{F,Q}
=
2\rho\delta+
\frac{3d-e}{4A^2\rho}\delta^2+
\bigO(\delta^3).
}
\]
For the \(P\)-axis,
\[
\boxed{
\kappa_{F,P}
=
2\rho\delta+
\frac{3g-f}{4A^2\rho}\delta^2+
\bigO(\delta^3).
}
\]

\subsection{Full sextic residual equilibrium}

At \(\delta=0\), write
\[
r=P/Q,
\]
\[
h_4(r)=A(1+r^4)+Br^2,
\]
\[
h_6(r)=d+er^2+fr^4+gr^6.
\]
Eliminating the radial variable from the stationary equations gives
\[
\boxed{
2h_4h_6'-3h_6h_4'=0.
}
\]
The physical root lies close to \(r=1\). The corresponding equilibrium is elliptic for the sixth-order Hamiltonian and provides the seed for the exact residual period-two cycles.
